\section{局部更好为何还不够}
\label{sec:09-Convex-Optimization/01-Optimization-Foundations/02}

同一份训练脚本跑两遍，只换了随机种子。第一遍验证损失停在 \(2.31\)，第二遍停在 \(2.05\)。两次的曲线都拉平了，两次的梯度范数都掉到 \(10^{-4}\) 量级，两次的日志都写着``已收敛''。

于是问题变成：第二次那个 \(2.05\)，是这个模型能达到的最好水平，还是又一个碰巧走不动的地方？如果再换十个种子，会不会出现 \(1.7\)？从算法内部看不出来。它掌握的全部信息是当前点的损失、梯度，或许还有一点曲率，而这些量描述的都是脚下这一小块地面。地面是平的，这件事千真万确；这块平地在整片地形里排第几，梯度不知道。

把话说精确。设在可行域 \(C\subseteq\R^n\) 上最小化 \(f\)。若存在半径 \(r>0\)，使

\[
f(x^\star)\le f(x)\quad\text{对一切 }x\in C\cap B(x^\star,r)
\]

成立，其中 \(B(x^\star,r)\) 是以 \(x^\star\) 为心、\(r\) 为半径的开球，就称 \(x^\star\) 为\textbf{局部最优}；若同一不等式对一切 \(x\in C\) 成立，才称它为\textbf{全局最优}。前者说的是``周围没有更好的''，后者说的是``所有可行点里没有更好的''。这两句话的差距，就是半径 \(r\) 与整个 \(C\) 的差距。

\begin{center}
\begin{tikzpicture}[x=1.6cm,y=0.9cm,font=\small,>=stealth]
  \node[font=\footnotesize,black!70] at (0.02,2.95) {非凸：两个坑，两个答案};
  \draw[black!75,thick] plot[domain=-1.55:1.6,samples=90] (\x,{(\x*\x-1)*(\x*\x-1)-0.3*\x});
  \draw[dashed,black!35] (-0.96,0.294) -- (-0.96,-0.72);
  \draw[dashed,black!35] (1.036,-0.305) -- (1.036,-0.72);
  \node[font=\footnotesize,black!70] at (-0.96,-1.0) {局部最优};
  \node[font=\footnotesize,black!70] at (1.036,-1.0) {全局最优};
  \fill[black!70] (-1.45,1.6505) circle (1.5pt);
  \fill[black!70] (-1.25,0.6914) circle (1.5pt);
  \fill[black!70] (-1.05,0.3255) circle (1.5pt);
  \fill[black!70] (-0.96,0.294) circle (1.5pt);
  \draw[->,black!55] (-1.45,1.6505) -- (-1.27,0.75);
  \draw[->,black!55] (-1.25,0.6914) -- (-1.07,0.38);
  \fill[black!70] (0.35,0.665) circle (1.5pt);
  \fill[black!70] (0.7,0.0501) circle (1.5pt);
  \fill[black!70] (0.95,-0.2755) circle (1.5pt);
  \fill[black!70] (1.036,-0.305) circle (1.5pt);
  \draw[->,black!55] (0.35,0.665) -- (0.66,0.11);
  \draw[->,black!55] (0.7,0.0501) -- (0.92,-0.23);
  \node[font=\footnotesize,black!60] at (-1.15,2.15) {初值 \(a\)};
  \node[font=\footnotesize,black!60] at (0.6,1.2) {初值 \(b\)};
  \begin{scope}[shift={(4.3,0)}]
    \node[font=\footnotesize,black!70] at (0.02,2.95) {凸：只有一个坑};
    \draw[black!75,thick] plot[domain=-1.55:1.6,samples=60] (\x,{0.9*\x*\x-0.3});
    \draw[dashed,black!35] (0,-0.3) -- (0,-0.72);
    \node[font=\footnotesize,black!70] at (0,-1.0) {局部最优 \(=\) 全局最优};
    \fill[black!70] (-1.4,1.464) circle (1.5pt);
    \fill[black!70] (-0.9,0.429) circle (1.5pt);
    \fill[black!70] (-0.4,-0.156) circle (1.5pt);
    \fill[black!70] (1.5,1.725) circle (1.5pt);
    \fill[black!70] (0.95,0.512) circle (1.5pt);
    \fill[black!70] (0.45,-0.118) circle (1.5pt);
    \fill[black!70] (0,-0.3) circle (1.5pt);
    \draw[->,black!55] (-1.4,1.464) -- (-0.98,0.55);
    \draw[->,black!55] (-0.9,0.429) -- (-0.48,-0.09);
    \draw[->,black!55] (1.5,1.725) -- (1.03,0.63);
    \draw[->,black!55] (0.95,0.512) -- (0.53,-0.05);
  \end{scope}
\end{tikzpicture}

\emph{左：同一个下降算法，两个初值，停在两个不同的谷底，而停下来的那一刻它分不出自己在哪一个。右：凸函数只有一个坑，走到走不动的地方，就是走到了最低点。}
\end{center}

图左边那两条轨迹用的是同一份代码、同一个步长规则，差别只在起点。它们各自的最后一步都满足``附近没有更好的方向''，也就都符合局部最优的定义，可两个终点的函数值差了 \(0.6\)。跑十次取最好的那一次，能提高撞上深坑的概率，却提高不了对结果的把握——你只是没找到更低的点，这和不存在更低的点是两回事。

\subsection{梯度为零只说明脚下是平的}

一元可微函数的内部极小点必有 \(f'(x)=0\)，这是微积分给的必要条件，而它离充分条件很远。\(f(x)=x^3\) 在原点导数为零，那里既不是极小也不是极大，只是拐弯前的一段平缓；\(f(x)=-x^2\) 在原点导数也为零，那里是全局最大。图左那条曲线有三个驻点，两个极小夹一个极大。驻点这个条件筛掉了大量的点，但它对``这块平地在整片地形中排第几''完全沉默。

维数一高，还会冒出\textbf{鞍点}：沿某些方向上升、沿另一些方向下降，梯度为零，Hessian 的特征值有正有负。高维非凸问题里鞍点通常远多于局部极小，算法在鞍点附近会长时间几乎不动，看上去像收敛。要把驻点再分类，得看二阶信息，见\hyperref[sec:09-Convex-Optimization/08-Nonconvex-and-Stochastic-Optimization/02]{``二阶结构区分局部极小与鞍点''}；即便分类清楚了，能拿到的保证也只到``接近某个驻点''为止，见\hyperref[sec:09-Convex-Optimization/08-Nonconvex-and-Stochastic-Optimization/01]{``下降引理在非凸地形中只保证接近驻点''}。

有约束时还要再退一步：最优点可能落在可行域边界上，那里梯度没有理由为零。在 \(x\ge1\) 上最小化 \(x^2\)，最优点是 \(x=1\)，导数是 \(2\)。可用的条件要换个说法——沿所有可行方向都无法一阶下降。把这句几何话写成一组可以逐条检验的等式与不等式，是后面\hyperref[sec:09-Convex-Optimization/05-Duality/03]{``KKT 把最优性写成一组方程''}的任务，这里先记下这个线索。

\subsection{凸性把两个词合成一个词}

现在给可行域和目标函数各加一条性质，看会发生什么。

设 \(C\) 是\textbf{凸集}：任取两个可行点，连接它们的整条线段仍在 \(C\) 内。设 \(f\) 在 \(C\) 上是\textbf{凸函数}：

\[
f((1-t)x+ty)\le(1-t)f(x)+tf(y),\qquad 0\le t\le1.
\]

几何上这说的是弦不低于图像，这句话会在后面两个单元里被反复展开。眼下只用它推一件事。

\begin{theorem}
凸函数在凸可行域上的每个局部最优点都是全局最优点。
\end{theorem}

证明的骨架只有三步。假设 \(x^\star\) 局部最优但不是全局最优，那么存在可行点 \(y\) 使 \(f(y)<f(x^\star)\)。在两点之间取 \(x_t=(1-t)x^\star+ty\)，其中 \(0<t\ll1\)。凸集保证 \(x_t\) 可行，凸函数保证

\[
f(x_t)\le(1-t)f(x^\star)+tf(y)<(1-t)f(x^\star)+tf(x^\star)=f(x^\star).
\]

而 \(t\) 足够小时 \(x_t\) 已经落进 \(x^\star\) 的邻域 \(B(x^\star,r)\)，于是邻域里出现了一个严格更低的可行点，与 \(x^\star\) 局部最优矛盾。

条件对一遍账。第一步用到 \(C\) 凸，否则线段可能穿出可行域，那个更低的点不可行；第二步用到 \(f\) 凸，且只用了定义里的那一个不等式；第三步用到局部最优的定义，靠 \(t\to0\) 把远处的 \(y\) 拉进半径 \(r\) 以内。全程没有用到可微、没有用到解的存在性、没有用到 \(C\) 有界，也没有假设 \(f\) 连续。用到的东西这么少，所以这条结论在不光滑的目标、带约束的问题和次梯度方法上同样成立。

回头看图左那条曲线，它在两个谷底之间先升后降，连接两个谷底的弦有很长一段落在图像下方，凸性不等式在这里被彻底违反。它能藏住更深的坑，靠的正是那道山脊：从浅谷出发必须先爬完一段上坡，才能到达深谷。凸性禁止的就是这种地形，于是``往任意更好的点走一小步就该立刻改善''这句话成立，局部证据被升级成了全局结论。

\subsection{凸不负责解存不存在、唯不唯一}

这条结论的边界要划清楚。它说的是``如果局部最优存在，它就是全局最优''，没说局部最优一定存在。\(f(x)=e^x\) 在整条实线上严格凸，下确界是 \(0\)，却没有任何一点取到它——这和上一节开区间上最小化 \(x\) 的例子是同一类事。存在性仍然要靠可行域的闭性与紧性，或者强制性条件（\(\norm{x}\to\infty\) 时 \(f(x)\to\infty\)）来提供。

唯一性也不在凸性的职责范围内。区间 \([-1,1]\) 上的常函数 \(f(x)=0\) 是凸的，每一点都是全局最优。要排除这种情况得加强假设：严格凸能保证``若最优解存在则至多一个''，论证方式和上面那三步一样——两个不同最优点的中点仍然可行，严格凸使中点的函数值严格低于二者共同的最优值，矛盾。再往前一步是强凸，它给出定量的曲率下界：离开最优点，函数值至少按二次增长抬起来。强凸不只带来唯一性，还直接决定算法的线性收敛速率和解对数据扰动的敏感程度，这三档条件的分工见\hyperref[sec:09-Convex-Optimization/03-Convex-Functions/03]{``曲率区分凸、严格凸与强凸''}。

\subsection{找到一个点和证明它最优是两件事}

非凸问题里常见的做法是多初值重启，几次都收敛到同一个点，信心就积攒了一些。这是经验，不是证明；没找到反例与反例不存在之间，隔着整个搜索空间。

凸问题能交出别的东西：一个\textbf{证书}。除了候选点本身，还能给出一个对所有可行点都成立的下界，两者的间隙有多小就写在纸面上。这个下界由对偶问题生成，见\hyperref[sec:09-Convex-Optimization/05-Duality/01]{``Lagrangian 怎样制造全局下界''}；间隙什么时候能闭合到零，是\hyperref[sec:09-Convex-Optimization/05-Duality/02]{``对偶间隙何时闭合''}的内容。证书的好处在于可以被独立核验——接手的人不必信任你的实现、你的步长、你的随机种子，只需要检查那几个不等式是否成立。

一份写着``损失 \(2.05\)，已收敛''的报告和一份写着``目标值 \(2.05\)，对偶下界 \(2.04\)，间隙不超过 \(0.5\%\)''的报告，交付的是两个等级的东西。现实中大量问题确实是非凸的，深度学习尤其如此，这时前一种报告已经是能拿到的最好结果；但只要模型还保得住凸性，就不该主动放弃后一种。

局部最优是站在原地能确认的事，全局最优是关于整个可行域的断言，一般问题里前者不蕴含后者，所以下降算法停下来时你并不知道自己在哪。凸性的全部价值，就是让这两件事合而为一——这一部分后面所有的内容，凸集、凸函数、对偶、KKT、内点法，都是在把这个条件辨认出来、构造出来、算出来。

那么就从最基本的问题开始：什么样的集合，能保证两个可行方案之间的整条线段都还可行。下一单元的\hyperref[sec:09-Convex-Optimization/02-Convex-Sets/01]{``线段、凸组合与凸包''}从这里接手。
